% ==============================================================================
% فصل هشتم: نتیجه‌گیری و کارهای آینده
% ==============================================================================

\chapter{نتیجه‌گیری و کارهای آینده}
\label{ch:conclusion}

\section{مقدمه}
در این فصل، دستاوردهای حاصل از طراحی، پیاده‌سازی و اعتبارسنجی عملیاتی سامانه تأمین داده ردیابی سخت‌افزاری جمع‌بندی شده و میزان تحقق تعهدات مصوب در پروپوزال اولیه مورد ارزیابی قرار می‌گیرد. همچنین با پیوند زدن یافته‌های فصول مختلف این رساله، مهم‌ترین درس‌آموخته‌های مهندسی مرور شده و نقشه راهی برای توسعه‌های آتی این پلتفرم در ابعاد دانشگاهی و کاربردهای پیشرفته صنعتی ترسیم می‌گردد.

\section{تطبیق خروجی‌ها با اهداف مصوب پروپوزال}
سامانه ارائه‌شده در این پایان‌نامه با هدف غلبه بر چالش دیرینه «اثر پروب» در استنتاج زمان اجرای بدترین شرایط (\lr{WCET}) توسعه یافت. جدول \ref{tab:proposal_compliance} میزان انطباق خروجی‌های پیاده‌سازی‌شده را با اهداف مصوب در پروپوزال رساله نشان می‌دهد:

\begin{table}[H]
\centering
\caption{جدول تطابق اهداف مصوب در پروپوزال با دستاوردهای عملی پیاده‌سازی‌شده}
\label{tab:proposal_compliance}
\begin{tabularx}{\textwidth}{|p{4.2cm}|X|c|}
\hline
\textbf{هدف مصوب در پروپوزال} & \textbf{شواهد عملی و خروجی پیاده‌سازی‌شده} & \textbf{وضعیت تحقق}\\
\hline
ردیابی غیرتهاجمی برنامه در سیستم نهفته & راه‌اندازی سخت‌افزاری \lr{ETMv4} و پورت \lr{TPIU} روی میکروکنترلر \lr{STM32H750} با اثبات ریاضی عدم تغییر رفتار زمانی برنامه (فصل ۵ و ۷) & محقق‌شده\\
\hline
استخراج دقیق زمان اجرا و چرخه‌های پردازشی & فعال‌سازی بسته‌های شمارنده چرخه (\lr{CCI}) و برچسب‌های زمانی سراسری (\lr{TS}) با تفکیک‌پذیری نانوثانیه در دکودر (فصل ۵ و ۶) & محقق‌شده\\
\hline
بازیابی جریان دستورات با تفکیک‌پذیری بالا & بازسازی بدون خطای جریان اجرای دستورات اسمبلی، پایش تصمیمات انشعاب و تشخیص مرزهای وقفه در سناریوهای آزمون تجربی (فصل ۶ و ۷) & محقق‌شده\\
\hline
توسعه محیط کاربری یکپارچه و اتوماسیون ابزاردقیق & طراحی و پیاده‌سازی افزونه اختصاصی \lr{CoreSight Trace Studio} در محیط \lr{VS Code} برای اتوماسیون کامل چرخه آزمون (فصل ۶) & محقق‌شده\\
\hline
تأمین بستر داده برای سامانه تست و تحلیل \lr{WCET} & تولید ساختار اسنپ‌شات استاندارد و لاگ‌های زمانی غنی جهت هم‌افزایی با پلتفرم آزمون خودکار محصول (فصل ۳ و ۷) & محقق‌شده\\
\hline
توسعه خط‌لوله با توجیه اقتصادی برجسته & جایگزینی تجهیزات تجاری چندهزار دلاری با ترکیب لاجیک‌آنالایزر دیجیتال، الگوریتم‌های پردازش سیگنال پایتون و برد کمکی (فصل ۴ و ۷) & محقق‌شده\\
\hline
\end{tabularx}
\end{table}

\section{سنتز و پیوند دستاوردهای فصول رساله}
پروژه حاضر پاسخی جامع به یک چالش مهندسی چندلایه بود که موفقیت آن مرهون هماهنگی دقیق میان لایه‌های سخت‌افزار، فرم‌ور، پردازش سیگنال و ابزارهای توسعه نرم‌افزار است:
\begin{enumerate}
    \item \textbf{تحقق نظارت کاملاً غیرتهاجمی بر پردازنده مدرن:}
    همان‌گونه که در فصول اول و دوم تبیین گردید، پردازنده‌های پیشرفته با معماری سوپراسکالر و خط‌لوله‌های عمیق (نظیر \lr{Cortex-M7})، به دلیل رفتارهای پیچیده پیش‌بینی انشعاب و حافظه‌های پنهان، امکان تحلیل ایستای دقیق \lr{WCET} را سلب می‌کنند و از سوی دیگر، ابزاردقیق نرم‌افزاری رفتار زمانی سیستم را به کلی دگرگون می‌سازد. در این پژوهش با راه‌اندازی ماژول سیلیکونی \lr{ETMv4} در فصل پنجم و ارائه مدل اثبات ریاضی در فصل هفتم، نشان داده شد که فواصل زمانی وقوع وقفه دوره‌ای ۱۰۰ میکروثانیه با انحراف زمانی کمتر از $0.1\ \mu\mathrm{s}$ ثبت می‌شود؛ این نتیجه اثبات قطعی حذف اثر پروب (\lr{Zero Probe Effect}) در این بستر است.

    \item \textbf{غلبه بر موانع فیزیکی ثبت سیگنال با نوآوری پردازشی:}
    در فصل چهارم، چالش یکپارچگی سیگنال‌های پورت پرسرعت \lr{TPIU} بر روی پایه‌های میکروکنترلر مورد واکاوی قرار گرفت. با طراحی راهبرد نمونه‌برداری ناهمگام ۱۰ برابری توسط لاجیک‌آنالایزر و توسعه اسکریپت تخصصی \lr{\texttt{TraceStreamProcessor.py}} در فصل ششم، مسائلی چون اعوجاج لبه‌ها، عدم تراز فاز و گلیچ‌های دیجیتال با الگوریتم‌های نرم‌افزاری برطرف گردید و خط‌لوله‌ای مستقل از تجهیزات گران‌قیمت آزمایشگاهی خلق شد.

    \item \textbf{همگام‌سازی بی‌نقص فرم‌ور تراشه با دکودر \lr{OpenCSD}:}
    یکی از مهم‌ترین درس‌آموخته‌های فنی این پروژه (که در فصول ۵، ۶ و ۷ تشریح شد)، وابستگی حیاتی موتورهای دکودر استاندارد به همخوانی دقیق مقادیر ثبّات‌های سخت‌افزاری پردازنده بود. کشف ریشه خطای \lr{NO\_SYNC} و تدوین دقیق فایل‌های اسنپ‌شات (\lr{\texttt{device1.ini}} و \lr{\texttt{device2.ini}})، پل مستحکمی میان جریان داده‌های خام و دکودر صنعتی ایجاد نمود.

    \item \textbf{ارتقای تجربه کاربری با توسعه افزونه اختصاصی \lr{CoreSight Trace Studio}:}
    برای جلوگیری از پیچیدگی‌های تنظیم دستی و تبدیل فرآیند دانشگاهی به یک محصول آماده برای صنعت، افزونه اختصاصی \lr{VS Code} در فصل ششم ارائه شد. این افزونه با پوشش خودکار مراحل تزریق درایور، پنجره‌گذاری کدهای مبدأ، ساخت اسنپ‌شات و مصورسازی شاخص‌های \lr{WCET} در قالب داشبورد گرافیکی، کل این فناوری پیچیده را به یک ابزار دم‌دستی و روان برای مهندسان سیستم‌های تعبیه‌شده تبدیل کرد.

    \item \textbf{هم‌افزایی ساختاریافته با سامانه آزمون خودکار محصول:}
    مطابق مرزبندی‌های معماری ارائه‌شده در فصل سوم و نتایج آزمون‌های فصل هفتم، بستر داده ارائه‌شده در این پایان‌نامه به عنوان موتور محرک زمانی عمل کرده و داده‌های استخراج‌شده از آن، خوراک حیاتی پلتفرم آزمون خودکار و ماک‌سازی حافظه‌محور (پروژه مکمل آقای عاملی) را فراهم می‌آورد؛ ترکیبی که در مقیاس کلی یک بستر بومی معادل سامانه‌های بین‌المللی چند ده‌هزار دلاری ایجاد نموده است.
\end{enumerate}

\section{چالش‌ها و محدودیت‌های مهندسی}
با وجود دستیابی کامل به اهداف مدنظر، این بستر در کاربردهای گسترده با چالش‌ها و محدودیت‌های زیر همراه است:
\begin{itemize}
    \item \textbf{وابستگی به اورسمپلینگ در فرکانس‌های کلاک بالا:}
    استفاده از لاجیک‌آنالایزر ناهمگام در فرکانس‌های ردیابی فراتر از ۵۰ تا ۱۰۰ مگاهرتز نیازمند نرخ‌های نمونه‌برداری چند گیگاهرتزی است که به محدودیت پهنای باند گذرگاه \lr{USB} و اشغال سریع بافر \lr{DRAM} می‌انجامد.
    \item \textbf{حجم بالای فایل‌های لاگ در پایش‌های طولانی‌مدت:}
    به دلیل صدور مداوم بسته‌های ردیابی و برچسب‌های زمانی، اجرای طولانی‌مدت برنامه‌ها حجم بالایی از داده‌های متنی را تولید می‌کند که نیازمند سازوکارهای فیلترینگ آنلاین و استریم مستقیم بر روی حافظه‌های پرسرعت است.
\end{itemize}

\section{پیشنهادها برای کارهای آینده}
جهت تکامل و گسترش کاربردهای این پژوهش، مسیرهای توسعه زیر پیشنهاد می‌گردد:
\begin{itemize}
    \item \textbf{توسعه بستر ثبت سخت‌افزاری سنکرون مبتنی بر \lr{FPGA}:}
    طراحی یک برد واسط مجهز به \lr{FPGA} ارزان‌قیمت که داده‌های پورت \lr{TPIU} را همگام با لبه‌های سیگنال \lr{TRACECLK} ثبت کرده و از طریق درگاه \lr{USB 3.0} یا گیگابیت اترنت به صورت استریم پیوسته ارسال نماید. این گام سامانه را قادر می‌سازد در حداکثر فرکانس نامی پردازنده (۴۸۰ مگاهرتز) بدون محدودیت بافر کار کند.
    \item \textbf{بصری‌سازی پیشرفته هیت‌مپ‌های زمانی در محیط ادیتور:}
    توسعه قابلیت‌های تکمیلی در افزونه \lr{CoreSight Trace Studio} جهت نمایش فلیم‌گراف (\lr{Flame Graph}) و هایلایت کردن خط به خط زمان اجرای توابع و مسیرهای بحرانی (\lr{Hotspots}) مستقیماً بر روی فایل‌های کد زبان \lr{C}.
    \item \textbf{اتصال برخط (\lr{Online Streaming}) به موتورهای تولید تست:}
    پیاده‌سازی یک لوله ارتباطی مبتنی بر سوکت شبکه (\lr{TCP Socket}) جهت انتقال آنی نتایج زمان‌بندی به الگوریتم‌های ژنتیک و موتورهای فازینگ آزمون، به نحوی که ورودی‌های تست به صورت خودکار به سمت تحریک طولانی‌ترین مسیرهای اجرایی هدایت شوند.
    \item \textbf{پشتیبانی از معماری‌های چند‌هسته‌ای و پردازنده‌های \lr{RISC-V}:}
    گسترش موتور پیش‌پردازش برای پشتیبانی از تراشه‌های چند‌هسته‌ای ناهمگن (\lr{Heterogeneous Dual-Core}) و تطبیق با استانداردهای ردیابی باز در خانواده پردازنده‌های \lr{RISC-V}.
\end{itemize}
